\section{Contamination information leakage and task lineage}\label{sec:C26:contamination}
Contamination is not a single binary property. A model may have seen the exact task, a near duplicate, the reference patch, a discussion of the solution, or a related repository version. A harness may leak hidden tests at runtime. Persistent memory can transfer answer-bearing information from previous evaluation episodes. Each channel requires a different control and supports a different inference about generalization.

Use task lineage to track duplicates and revisions. A later benchmark release can contain transformed versions of earlier public tasks. Deduplication based only on exact text can miss equivalent code changes or renamed entities. Semantic similarity is useful for screening but can also flag legitimately related tasks. Report the screening procedure and unresolved cases rather than claim complete decontamination from one similarity threshold.

Temporal holdouts reduce some exposure but do not guarantee novelty. A task first published after model training may depend on a public issue or patch that predates the cutoff. Conversely, a benchmark revision date does not make an old task new. Record first public availability of the task and solution separately when evidence permits. For proprietary models, unknown training data limits what can be established.

Runtime leakage should be tested directly. Keep evaluator credentials outside the agent environment, separate reference artifacts, and inspect tool access. A browser agent can inadvertently observe a success marker or database field intended only for grading. A coding agent can discover test files copied into the wrong layer. These are implementation defects that can be fixed without speculative claims about model memorization.

Memory experiments require explicit contamination boundaries. If one arm carries experiences across tasks while another resets, the information difference is the treatment and must be reported. If a learned skill contains a benchmark-specific answer rather than a reusable procedure, transfer gains can reflect memorization. Holdouts should test both new instances and changed task structure, with source and experience provenance available for audit.

\subsection{Interference through persistent state}\label{sec:C26:stateinterference}
A stateful benchmark generally does not admit independent per-task potential outcomes. Let $X_{1:t}$ be the task prefix, $M_t^S$ the memory generated by system $S$ on that prefix, and $Y_t(S,M_t^S,X_t)$ the next outcome. Changing the system at an early task changes later memory and therefore later outcomes even when the task list is fixed. A mean over tasks estimates performance of the complete stream policy under its order and feedback rules; it does not identify an isolated memory effect at every task.

Two experiments answer different questions. A frozen-prefix experiment initializes both arms from the same admissible state and varies the next decision mechanism. A full-stream experiment initializes a common starting state and lets each arm construct its own subsequent memory. The former isolates a conditional local effect; the latter measures accumulated adaptation and error propagation. Resetting memory after every task removes the longitudinal treatment, while copying the stronger arm's evolved bank into both conditions changes the treatment again. Report the chosen intervention and cluster uncertainty by independently generated stream. Chapter~\ref{ch:C27} develops the corresponding counterfactual controls and contamination boundaries.
